\subsection{Codimension-two pencils and the trace image}
\paragraph{J6P (Proved with the stated smooth-axis hypothesis).}
The middle assertion at 44:4 reads ``le facteur fixe de $P^i(Y)$
est identique à l'image de $P^i(X)$''. Here $\operatorname{Ab}^2$ means the universal regular algebraic
representative, not an unidentified manuscript quotient.
Let $k$ be algebraically closed, let $X/k$ be smooth projective of
dimension $n\geq4$, and let a hyperplane pencil have a smooth
codimension-two base locus $B$. Assume its generic member $Y/K$,
$K=k(t)$, is smooth. Write $j:Y\hookrightarrow X_K$ and
$A=\operatorname{Ab}^2_{Y/K}$. Then
\begin{equation}\label{eq:J6P-trace}
 \operatorname{im}\bigl(j^*:\operatorname{Ab}^2_{X/k,K}\longrightarrow A\bigr)
 =\operatorname{im}\bigl(\tau:T_K\longrightarrow A\bigr),
 \qquad T=\operatorname{Tr}_{K/k}(A),
\end{equation}
as abelian subvarieties of $A$. If $n\geq5$, both images are $A$,
so $j^*$ is surjective. No standard conjecture or torsion
injectivity assumption is used for these two conclusions.
The smoothness of the pencil axis is an additional hypothesis:
it is not inferred from the manuscript's integral-fibre condition.

The representative and its separable base-change compatibility
are supplied by \cite[Theorems 6.1 and 5.10, pp.~25--27]{DGJ_ACMVFunctor}.
The trace exists by \cite[Definition 6.1 and Theorem 6.2, p.~20]{DGJ_Conrad2006}.
We use its image, since its kernel can be infinitesimal in positive
characteristic; see \cite[Theorem 6.12, pp.~22--23]{DGJ_Conrad2006}.

\emph{Step 1: spreading a trace family.}
Set $\widetilde X=\operatorname{Bl}_B X$, with blowdown $p$,
exceptional divisor $e:E\hookrightarrow\widetilde X$, and
$q:E\to B$. The pencil gives $f:\widetilde X\to\mathbb P^1$.
Its generic fibre is $Y$: blowing up a member along its Cartier
divisor $B$ changes nothing. Both $\widetilde X$ and $B$ are
smooth projective over $k$.

Choose a miniversal cycle
$\Gamma\in\operatorname{CH}^2(A\times_KY)$ whose induced map
under the universal regular homomorphism $\Phi_Y$ is $[r]_A$,
where $r>0$. This exists over $K$ itself, including when $K$ is
imperfect, by \cite[Lemma 4.7, pp.~18--19]{DGJ_ACMVFunctor}.
Pull it back along $\tau\times1$ to
$u\in\operatorname{CH}^2(T_K\times_KY)$. Normalize at the zero
of $T$ if necessary. Then
\[
 \Phi_Y(u_t-u_0)=r\tau(t).
\]
Take closures of cycle representatives to obtain
$Z\in\operatorname{CH}^2(T\times_k\widetilde X)$ restricting
to $u$ over the generic point of $\mathbb P^1$. Refined fibres
$Z_t-Z_0$ are algebraically trivial and give a regular
homomorphism $v:T\to\operatorname{Ab}^2_{\widetilde X/k}$.
The generic restriction $\delta$ therefore satisfies
\begin{equation}\label{eq:J6P-spread}
 \delta\circ v_K=r\tau,
 \qquad
 \delta:\operatorname{Ab}^2_{\widetilde X/k,K}\longrightarrow A.
\end{equation}
This equality holds as morphisms: it is the equality of regular
maps induced by the generic restriction of the same Chow class.

\emph{Step 2: the exceptional summand.}
The codimension-two blowup formula is
\[
 \operatorname{CH}^2(\widetilde X)
 =p^*\operatorname{CH}^2(X)\oplus e_*q^*\operatorname{CH}^1(B).
\]
Its inverse is $(p_*,-q_*e^*)$. Indeed,
$p_*p^*=1$, $p_*e_*q^*=0$, $q_*e^*p^*=0$, and
$q_*e^*e_*q^*=-1$, the last identity following from
$N_{E/\widetilde X}=\mathcal O_E(-1)$ and
$q_*c_1(\mathcal O_E(1))=1$. The general formula is opened in
\cite[Section 2.3, Theorem 2.8(3)]{DGJ_Jiang2023}, which states
Fulton's blowup formula. All these maps preserve algebraic
equivalence and commute with the refined fibres in a smooth
parameter space. Consequently the decomposition is compatible
with regular homomorphisms and gives an actual isomorphism
\begin{equation}\label{eq:J6P-blowup}
 \operatorname{Ab}^2_{\widetilde X/k}
 \simeq\operatorname{Ab}^2_{X/k}\times\operatorname{Pic}^{\rm var}_{B/k}.
\end{equation}
One can also obtain this directly from the universal properties:
the displayed inverse Chow operations give inverse homomorphisms
of the three universal regular representatives.

Let $b:B_K\hookrightarrow Y$. Under
\eqref{eq:J6P-blowup}, $\delta$ is the sum of $j^*$ and the
homomorphism induced by $b_*$ on algebraically trivial divisors.
There is no excess term here: $E$ and the smooth generic fibre
meet in the Cartier divisor $B_K$, so refined base change gives
$\delta e_*q^*=b_*$.

Choose a smooth member $Y_0/k$ of the pencil; one exists since
the smooth locus of the proper pencil contains its generic point.
Applying weak Lefschetz to $X\supset Y_0\supset B$ shows that
restriction
\[
 a:\operatorname{Pic}^{\rm var}_{X/k}
       \longrightarrow\operatorname{Pic}^{\rm var}_{B/k}
\]
is an isogeny. Indeed, $n\geq4$ makes both restrictions
isomorphisms on $H^1(-,\mathbb Q_\ell(1))$ and hence on the
rational Tate modules of these Picard varieties. Write $H$ for
the chosen hyperplane class. The projection formula gives,
for every algebraically trivial line bundle $L$ on $X$,
\[
 b_*(L|_{B_K})=H|_Y\cdot L|_Y=j^*(H\cdot L).
\]
Thus, as homomorphisms of abelian varieties,
$b_*a_K=j^*\beta_K$, where
$\beta:\operatorname{Pic}^{\rm var}_{X/k}\to
\operatorname{Ab}^2_{X/k}$ is induced by $L\mapsto H\cdot L$.
Since $a$ is surjective, the exceptional summand has image in
$\operatorname{im}(j^*)$. Therefore
$\operatorname{im}(\delta)=\operatorname{im}(j^*)$.

Equation \eqref{eq:J6P-spread} and the surjectivity of $[r]_T$
now imply $\operatorname{im}(\tau)\subseteq\operatorname{im}(j^*)$.
Conversely, the map $j^*$ has constant source; the universal
property of $T$ factors it through $\tau$. This proves
\eqref{eq:J6P-trace}.

\emph{Step 3: why the trace fills the stable-range target.}
Assume $n\geq5$. Choose a prime $\ell$ different from the
characteristic and coprime to the integer of a miniversal cycle
for $A$. By \cite[Proposition 8.2, p.~31]{DGJ_ACMVFunctor}, there
is a Galois-equivariant injection
\[
 T_\ell A\hookrightarrow H^3(Y_{K^{\rm sep}},\mathbb Z_\ell(2))_{
 \rm torsion\,free}.
\]
Weak Lefschetz identifies its rational target with
$H^3(X_{K^{\rm sep}},\mathbb Q_\ell(2))$. This Galois
representation is constant because $X$ is defined over the
algebraically closed field $k$. Thus $T_\ell A$ has trivial
Galois action. All $\ell$-power torsion points of $A$ are
$K$-rational, and their union is a divisible group.

By Lang--N\'eron,
$A(K)/\tau(T(k))$ is finitely generated
\cite[Theorem 7.1, p.~23]{DGJ_Conrad2006}. A divisible group
has zero image in a finitely generated abelian group, so
$A(K)[\ell^\infty]\subseteq\tau(T(k))$. Prime-to-characteristic
power torsion is Zariski dense in an abelian variety. For
completeness, if a proper abelian subvariety contained all this
torsion, its positive-dimensional abelian quotient would have
no $\ell$-power torsion; this is impossible since torsion lifts
through a surjective homomorphism of abelian varieties over an
algebraic closure. Hence $\operatorname{im}(\tau)=A$, as claimed.

\emph{The finite-kernel boundary.}
For $n\geq4$, the same $j^*$ has finite kernel if the universal
map for $Y$ is injective on $\ell$-primary cycle torsion for a
prime $\ell$ avoiding the characteristic and the integer of a
chosen miniversal cycle for $X$. Here is the full implication.
If the reduced identity component $C$ of its kernel were
positive-dimensional, restrict a miniversal cycle $\Gamma_X$
to $C$. The map
\[
 T_\ell C\longrightarrow T_\ell A^2(X_{K^{\rm sep}})
 \longrightarrow H^3(X_{K^{\rm sep}},\mathbb Z_\ell(2))_{
 \rm torsion\,free}
\]
is injective by Proposition 8.2. Its pullback cycle family on
$Y$ is killed by $\Phi_Y$ since the representative map is
$[r]j^*|_C=0$. The stated torsion injectivity then makes its
Tate module of cycles zero. Functoriality of the Bloch map
\cite[Section 8.1, pp.~30--31]{DGJ_ACMVFunctor} makes the
displayed cohomology classes restrict to zero on $Y$. But weak
Lefschetz is injective on $H^3(X)\to H^3(Y)$ for every $n\geq4$,
a contradiction after tensoring with $\mathbb Q_\ell$.

This torsion hypothesis is known in characteristic zero by
Murre's theorem, as recorded in
\cite[Remark 8.6, p.~33]{DGJ_ACMVFunctor}. It therefore gives,
in characteristic zero and with the smooth-axis hypothesis,
an isogeny in the stable range and an isogeny onto the trace
image when $n=4$. Proposition 8.2 alone does not remove the
hypothesis: its injection is constructed from a miniversal
cycle and does not assert that $\Phi_Y$ is injective on torsion.
The same remark explicitly leaves that injectivity unsettled
in general positive characteristic.

\paragraph{J6Z (The highest nontrivial target codimension).}
Let $X/k$ be smooth projective of dimension $n\geq3$ and let
$j:Y\hookrightarrow X$ be any smooth hyperplane section. Then
\[
 j^*:P^{n-1}_{\rm inc}(X)\longrightarrow
       P^{n-1}_{\rm inc}(Y)=\operatorname{Alb}(Y)
\]
is surjective in every characteristic. The analogous assertion
holds for universal regular representatives whenever the source
exists; in particular it holds for
$\operatorname{Ab}^2(X)\to\operatorname{Alb}(Y)$ on a threefold.

To prove it, restriction
$\operatorname{Pic}^{\rm var}(X)\to\operatorname{Pic}^{\rm var}(Y)$
is an isogeny by weak Lefschetz in degree one. Intersection by
$H_X^{n-2}$ defines a homomorphism from the former Picard
variety to $P^{n-1}_{\rm inc}(X)$, and its restriction to $Y$
is intersection by $H_Y^{n-2}$ followed by the Albanese map.
The induced map
\[
 \operatorname{Pic}^{\rm var}(Y)\longrightarrow\operatorname{Alb}(Y)
\]
is an isogeny: on rational Tate modules, the codimension-one
and zero-cycle Bloch isomorphisms identify it with
\[
 L_Y^{n-2}:H^1(Y,\mathbb Q_\ell(1))
       \xrightarrow{\ \sim\ }H^{2n-3}(Y,\mathbb Q_\ell(n-1)),
\]
which is hard Lefschetz. See
\cite[Section 8.1, pp.~30--31]{DGJ_ACMVFunctor} for those
Bloch isomorphisms and
\cite[Theorem 4.1.1 and (4.1.6), pp.~217--219]{DGJ_Deligne1980}
for hard and weak Lefschetz. A surjective homomorphism thus
factors through $j^*$, proving surjectivity. In this top target
codimension, the common incidence quotient is the Albanese:
the Poincar\'e family of divisors on $Y$ detects exactly its
kernel. This argument uses neither an inverse Lefschetz
correspondence nor the torsion standard assumption.

\emph{Exact scope.}
The fixed-part result proves the modern codimension-two
interpretation of the middle clause for pencils with smooth
axis. It does not identify that interpretation with the
manuscript's undefined equivalence relation, cover singular
pencil axes, prove positive-characteristic finite kernel, or
settle the higher-codimension stable and middle assertions.
